%! HV lualatex --synctex=1
%$ cllualatex --synctex=1 $texname
\documentclass{simplman}

\usepackage{inlinegraphicx}
\usepackage{ninecolors}
\usepackage[shortverb=false,private,backcolor=gray9!30,framecolor=gray5!50]{srcandtext}

\usepackage{kswdash}
\usepackage[localenv,parentype=moreparens]{kskorsymb}
\let\ksemph\respectdotemph
\usepackage{ksvqt}
\newcommand*\tblrdoumi{\<3\textsf{T/B/L/R 도우미}>}

\newcommand*\env[1]{\<2\textit{#1}>}
\let\pkg\stp
\let\opt\stopt

\newcommand*\guielem[1]{\colorbox{blue8!20}{\sffamily[#1]}}

\ExplSyntaxOn

\NewDocumentEnvironment {widetext} {}
{
	\begin{adjustwidth}{-.5\spinemargin}{0pt}
}
{
	\end{adjustwidth}
}

\NewDocumentCommand \applesymbolfont { m }
{
	\fontspec{Apple~Symbols.ttf}[Scale=#1]
}

\ExplSyntaxOff

\usepackage{multirow}

\maxtocdepth{subsection}

\begin{document}

\begin{titlingpage}
\calccentering{\unitlength}
\begin{adjustwidth}{\unitlength}{-\unitlength}
\begin{center}
\inlinegraphics[scale=5]{tblricon}
\huge \tblrdoumi{} 3.9 사용 설명서\\
\large Nova de Hi\quad 2026/04/02
\end{center}

\vspace{-\baselineskip}

\tableofcontents*
\end{adjustwidth}

\end{titlingpage}


\chapter{표(tabular)에 대한 관점}

\sidepar{%
이 글에서 `tabular'라고 하면 표 자체를 가리킨다. 반면 \env{tabular}는 \LaTeX 의 \env{tabular} 환경을 의미한다.
}

표는 다음과 같이 이루어져 있고 이 각각은 다른 것과 구분되는 요소이다.

\begin{enumerate}[\quad(1)]\oblivoirlist
\item 표의 본체(body). array. \verb|&|와 \verb|\\|로 이루어지는 이차원 배열(array) 구조로 나열된 텍스트.
\item 괘선. 본체의 배열 구조를 둘러싸는 선.
\item 행(row)\.열(column)\.셀(cell)의 속성 설정
\end{enumerate}

``구분된다''는 표현의 의미는 서로 독립적이라는 말이다. 표의 본체\ksparensmall{코드}에 아무런 영향을 주지 않고 괘선의 굵기나 색상을 달리할 수 있다. 표의 본체에 손대지 않고 표의 정렬상태나 폰트, 색상을 조절할 수 있(어야 한)다.

지난 40년 \LaTeX\ 역사에서 표 그리기는 \env{tabular}에 의존해왔다. 
그 결과 세 가지 \ksemph{구분되어야 할} 요소들이 뒤엉킨 채로 표를 작성하는 데 익숙해지고 말았다.
복잡한 것을 추가로 구현하게 될수록 소스는 도저히 읽을 수 없는 것이 되어갔다.

우리는 가능한 한 이 세 요소를 구분하여 다루고자 한다. 
그리고 이것을 가능하게 하는 것이 \stp{tabularray}의 \env{tblr} 환경이다.
\tblrdoumi 는 \env{tblr} 환경의 표를 생성하게 하는 GUI 유틸리티이고, 표를 바라보는
``새 관점''에 입각해 있다.

\chapter{최초 실행 상태}

\begin{center}
\includegraphics[scale=.5]{screenshot008}
\end{center}

\begin{itemize}
\item 최상단 왼쪽은 (편집불가능한) \textcolor{green5}{colspec 미리보기}이다.
\begin{itemize}
\item colspec 문자열이 \texttt{X[c,m]}의 3번 반복으로 되어 있는 것은 최초 선택되어 있는 colspec mode가 \opt{linewidth}이기 때문이다.
\end{itemize}
\item 최상단 오른쪽에는 버튼이 세 개 있다. 초록색은 ``Excel 파일 불러오기'', 파란색은 ``메뉴 열기'', 그리고 [설정 패널]을 열거나 닫는 토글 버튼.
\item \textcolor{green5}{테이블 영역}에는 2$\times$3 테이블이 예시되어 있다.
\begin{itemize}
\item 1번 행에 마크가 붙어 있는 이유는 최초 예시에 \texttt{row\{odd\}=\{bg=gray!10\}}이라는 행 설정이 되어 있기 때문이다. 그러므로 만약 3번 행을 이 상태에서 만든다면 3, 5번등 홀수행에 이 마크가 모두 붙을 것이다.
\end{itemize}
\item 괘선이 두드러지게 그어져 있는데 이것은 예시의 Lines 설정 때문이다. 괘선을 다루는 장에서 상세히 설명한다.
\item \textcolor{green5}{설정 패널}의 [테이블 환경]은 \opt{tblr}, \opt{longtblr}, \opt{talltblr} 중에 고르는 콤보박스로 되어 있다.
\item {} [기본 설정] 탭과 [long/tall 설정] 탭이 있는데 [long/tall 설정] 탭은 비활성화되어 있다. 이것은 위의 [테이블 환경]을 \opt{longtblr}이나 \opt{talltblr}로 하면 열 수 있는 상태가 된다.
\item {} [Tab Size]는 스핀박스를 이용하여 행과 열의 수를 조절할 수 있고 [적용] 버튼이 달려 있다.
\item {} [Colspec Mode] 다음에 3개의 라디오 버튼이 있는데 최초에는 \opt{linewidth}가 선택되어 있을 것이다.
\item 그 오른쪽에 있는 [colspec 편집] 버튼은 \opt{manual} 모드와 깊은 관계가 있다.
\item 아래 [manual colspec] 라인에디터에는 아무 내용도 없으며 이것은 [colspec 편집]에서 다루는 ``manual colspec''의 내용 확인 창이다. 즉 이 영역을 직접 편집할 수 없고 [colspec 편집] 버튼을 통해서 한다.
\item 그 아래 네 개의 설정 편집기들이 달려 있다. 이 설정 편집기는 각각이 독자적인 에디터이며 따로 떼어내거나 탭으로 겹치거나 할 수 있다. 
\item 하단에는 \textcolor{green5}{셀 편집 영역}과 \textcolor{green5}{Latex 코드 미리보기}가 붙어 있다.
\end{itemize}

이 가운데 크기가 고정되어 있는 것은 \textcolor{green5}{설정 패널}이다. 이것을 가리고 창의 가장자리를 이동하여 창 크기를 줄이거나 경계선 부근을 끌어서 각 영역이 차지하는 부분을 조절할 수 있다.

\begin{Chamgo}
최초 실행은 그 자체가 하나의 작은 예제가 되게 되어 있으나 새로운 표 작성을 시작할 때는 이 미리 주어진 설정이 방해가 될 수 있다. 초기 괘선\.행\.열\.셀 설정을 깨끗이 지우고 작업을 시작하는 것이 좋으며 이를 위해서 (하나씩 찾아서 지워도 되지만) \hcrkey{F9} 단축키를 이용하라. 이 단축키는 포커스가 테이블 영역에 있을 때 동작한다. 즉 최초 실행 상태에서 \hcrkey{ESC}, \hcrkey{F9}를 차례로 누르면 된다.
\end{Chamgo}

시스템 테마가 dark일 때는 어두운 색상을 반영한다. 그러나 프로그램 자체에 테마 기능을 가지고 있지는 않다.

\hcrkey{F1}을 누르면 단축키를 요약해서 보여준다.

\chapter{표의 본체}

표를 작성할 때 맨 먼저 해야하는 일은 표의 본체 array 구조를 만드는 일이다.
\LaTeX 에서 이 구조는 \env{array}나 \env{matrix}와 완전히 같다.
그리고 이 구조를 만드는 동안은 \ksemph{괘선}이나 \ksemph{폰트, 정렬, 색상}에 대해서는 잊어버리는 것이 좋다.

\env{tabular}의 경우에는 이것이 불가능했다. 다음 예를 보자. 이 예제는 \href{https://latex-tutorial.com/tutorials/tables/}{\ttfamily latex-tutorial.com}에서 가져왔다.

\begin{widetext}
\setrighthandwidth{.4\linewidth}
\begin{exampleside}
\begin{tabular}{l|c|r}
\textbf{Value 1} & \textbf{Value 2} & \textbf{Value 3}\\
$\alpha$ & $\beta$ & $\gamma$ \\ \hline
\multicolumn{2}{c|}{12} & a\\ \hline
2 & 10.1 & b\\
3 & 23.113231 & c\\
4 & 25.113231 & d\\
\end{tabular}
\end{exampleside}
\end{widetext}

여기서 다음과 같은 점을 생각해볼 수 있다.
\begin{itemize} \oblivoirlist
\item \verb|\textbf|는 표의 내용이 아니다. 그것은 ``장식적인'' 요소일 뿐이다.
\item \verb|\hline|은 괘선이다. 이 역시 표의 내용에 해당하지 않는다.
\item \verb|\multicolumn|은 셀의 병합을 위하여 쓴 것인데 이 명령이 쓰인 행은 다른 행과 비교하여 \verb|&|의 개수가 줄어든다. 이렇게 설계한 이유는 아마도 코딩의 편의를 위한 것이었겠지만 \verb|&|의 수로만 보면 일관성이 없어졌다.
\end{itemize}

\env{tabular}는 이와 같이 표의 내용과 그 스타일링 요소를 뒤섞어 입력해야 했던 것이다.
이것들을 모두 들어내고 오직 표의 본체만을 적어보면 다음처럼 된다.

\begin{exampleonly}
Value 1  & Value 2 & Value 3  \\
$\alpha$ & $\beta$ & $\gamma$ \\
12 &           & a\\
2  & 10.1      & b\\
3  & 23.113231 & c\\
4  & 25.113231 & d\\
\end{exampleonly}

\section{입력}

\tblrdoumi 를 실행하여 표의 본체를 만드는 데까지 진행해보겠다.

\begin{center}
\includegraphics[scale=.5]{screenshot001}
\end{center}

\begin{enumerate}[설명 1]\oblivoirlist
\item 처음 실행한 상태는 \guielem{설정 패널}이 열려 있는 상태이다. 상단의 검은색 버튼을 눌러서 열거나 닫을 수 있다.
\item 화면의 오른쪽 하단에 현재 포커스가 있는 셀의 위치가 표시된다. 맨처음 실행하면 \texttt{(0,0)}이라고 나오는데 이것은 아직 \guielem{테이블 영역}에 진입하지 않았다는 뜻이다. 셀을 클릭하거나 \hcrkey{ESC}를 누르면 셀을 선택하는 상태가 된다.
\item 셀에 문자를 입력하기 시작하면 셀 편집 상태가 된다. 셀/열/행에 대하여 일정한 편집적 조작을 가하기 위해서는 셀 선택 상태이어야 하고 셀 편집 상태에서 빠져나와야 한다. 텍스트 입력 후 \hcrkey{Enter}나 TAB(\hcrkey{Tab})을 누르면 셀 편집 내용이 확정되어 셀 선택 상태가 된다. 셀 편집 상태에서 \hcrkey{ESC}는 입력 내용을 취소하는 것이 되므로 주의.
\item 셀에 입력해야 할 문자열이 길고 복잡한 때는 좌측 하단의 \guielem{셀 편집 영역}을 이용하여 입력할 수 있다. 이 영역에는 ``선택된'' 셀의 내용이 표시되는데, 셀 편집 상태에서는 아직 그 내용이 확정되지 않았기 때문에 나타나지 않는다. 반면 \guielem{셀 편집 영역}에 입력하는 문자열은 테이블 영역의 셀에 즉시 나타난다.
\item 셀에서 \guielem{셀 편집 영역}으로 바로 이동하는 단축키는 \hcrkeys[+]{Ctrl, Insert} 또는 \hcrkeys[+]{Ctrl, E}(\uline{E}dit)이다. 셀 편집 영역에서 편집이 끝나면 \hcrkey{ESC}를 눌러서 \guielem{테이블 영역}으로 복귀한다.
\item 셀 선택 상태에서는 각 셀이나 행\.열을 선택하여 복사\.잘라내기\.붙여넣기할 수 있다.
\end{enumerate}

행\.열\.셀을 선택하는 방법은 다음과 같다.
\begin{itemize}
\item 클릭: 한 개의 셀을 선택한다.
\item \hcrkey{Ctrl}+클릭: 개별 셀들을 하나씩 선택한다.
\item \hcrkey{Shift}+이동키: 인접한 셀들을 선택한다.
\item 마우스를 클릭한 후 드래그하여 여러 셀을 선택할 수 있다.
\item 상단 열번호가 나열된 헤더 부분을 클릭하거나 드래그하거나 \hcrkey{Shift}, \hcrkey{Ctrl} 키와 함께 열 전체를 선택할 수 있다.
\item 좌측 행번호가 나열된 헤더 부분을 클릭하거나 드래그하거나 \hcrkey{Shift}, \hcrkey{Ctrl} 키와 함께 행 전체를 선택할 수 있다.
\item 헤더의 최상단 왼쪽 빈 공간을 누르면 모든 셀이 포함되는 테이블 전체를 선택할 수 있다.
\end{itemize}

\begin{Chamgo}
\begin{itemize}\oblivoirlist
\item 이미 텍스트가 있는 셀에 대하여 포커스를 두고 텍스트 입력을 시작하면 새로운 내용으로 대체할 수 있다. 기존 텍스트를 그대로 두고 편집 상태로 들어가려면 더블클릭한다.
\item 표의 본체만 입력하기 위해서는 \textcolor{blue4}{[설정 패널]}을 닫은 상태에서가 편하다.
\item 열(columns) 전체를 선택하여 header의 열 가장자리를 끌면 선택한 전체 열이 같은 폭으로 늘어난다. (하나의 열에 대해서도 열의 폭을 늘리거나 줄일 수 있다.)
\item \textcolor{blue4}{셀 에디터 영역}은 긴 텍스트를 입력할 때 쓸 수 있고 현재 테이블에서 일부 가려진 셀 텍스트를 확인하기 위해서도 쓸 수 있다.
\item 셀 선택 상태에서 행\.열\.셀을 선택한 후에 복사하여 Excel이나 구글 스프레드시트에 붙여넣을 수 있다.
\item 최초 실행 상태에서 텍스트를 입력하고 확정하면 텍스트가 셀의 가운데로 온다. 그 이유는 colspec mode가 \textcolor{blue4}{\ttfamily linewidth}이기 때문이다. 모두 왼쪽으로 오도록 입력하고 싶다면 (셀 선택 상태에서) \hcrkey{F9}를 누르거나, 또는 \textcolor{blue4}{\ttfamily natural}로 모드를 바꾸면 되는데 상세한 사항은 Colspec 관련 절에서 설명한다. 표 본체 입력시에는 텍스트의 정렬 상태에는 신경쓰지 않는 것이 좋다.
\end{itemize}
\end{Chamgo}

표의 본체 입력이 진행되는 동안 우하단의 \guielem{LaTeX 코드 미리보기} 영역에 생성될 코드가 나타난다.

\section{찾기와 바꾸기, 위치 찾기}

텍스트를 입력하는 상황에서는 항상 찾기와 바꾸기가 필요하다. 

만약 포커스가 \guielem{테이블 영역}에 있다면 \hcrkeys[+]{Ctrl, F}로 찾기를 실행할 수 있다. 주어지는 텍스트가 포함된 모든 셀을 차례로 찾는데, \hcrkey{F3}으로 다음 위치, \hcrkeys[+]{Shift,F3}으로 이전 위치로 갈 수 있다.

바꾸기는 ``일괄 치환'' 기능만 있는데 \hcrkeys[+]{Ctrl, H} 단축키로 동작한다. (맥에서는 Cmd+H 단축키가 시스템에 예약되어 있으므로 우클릭 메뉴에서 ``모두 바꾸기''를 선택하는 것이 좋다.)

\begin{Chamgo}
텍스트를 포함한 셀로 찾아가기와 테이블 내의 텍스트 치환 동작은 포커스가 \guielem{테이블 영역}에 있을 때만 유효하다. 다른 편집 영역에 포커스가 가 있으면 이 단축키는 그 포커스 내부의 편집 기능으로 동작한다.
\end{Chamgo}

표를 편집하다가 보면 특정 행\.열\.셀 위치로 바로 가야 할 때가 있다. 이를 위한 단축키가 몇 개 마련되어 있다.
\begin{description}
\item [{\hcrkeys[+]{Ctrl,G}}] 숫자 입력을 받아들여서 해당 행으로 이동한다.
\item [{\hcrkeys[+]{Shift, Ctrl,G}}] 숫자 입력을 받아들여서 현재 포커스가 위치한 행의 주어지는 열로 이동한다.
\item [{\hcrkeys[+]{Ctrl,Alt,G}}] 두 개의 숫자를 쉼표로 구분하여 받아들여서 해당 셀로 이동한다. 예를 들면 \texttt{3,5}라고 입력하면 (3,5) 셀로 이동한다. 일부 기기에서 이 단축키를 다른 앱이 사용하는 경우가 많아서 \hcrkeys[+]{Shift,Alt,G}로도 같은 기능을 하도록 해두었다. (참고. Mac에서는 Option+Shift+G({\applesymbolfont{1.5}⌥⇧}G)이다.)
\end{description}

\noindent 
이밖에 \hcrkey{home}, \hcrkey{end}, \hcrkeys[+]{Ctrl, Home}, \hcrkeys[+]{Ctrl, End}는 각각 행의 첫 셀로, 행의 마지막 셀로, 열의 첫 셀로, 열의 마지막 셀로 포커스를 이동시킨다.

\section{표의 크기}

행과 열의 수를 \guielem{설정 패널}에서 정할 수 있다. 
그리고 \guielem{테이블 영역}의 셀에 포커스가 있을 때 다음 단축키를 이용하여 행\.열을
삭제하거나 추가할 수 있다. Table Size의 행/열 수는 이에 맞춰 변화한다.

\begin{itemize}
\item \hcrkey{F7}. 최상단 행($=$ 1번행)에 포커스가 있으면 오른쪽에 새 열을 만든다.
\item \hcrkey{F8}. 최상단 행($=$ 1번행)에 포커스가 있으면 현재 열을 삭제한다.
\item \hcrkeys[+]{Ctrl, Enter}. 행의 마지막 셀에 포커스가 있으면 아래에 새 행을 만든다.
\item \hcrkeys[+]{Ctrl, Delete}. 현재 행을 삭제한다.
\end{itemize}

\section{들여오기}

어차피 표의 본체를 입력하는 것만이 목표라면 굳이 이 앱에서 하지 않아도 될 것이다. 즐겨 다루는 Microsoft Excel 또는 구글 스프레드시트에서 표의 ``본체''를 작성하는 것이 더 쉬울 때도 있다.
\tblrdoumi 는 이러한 상황에 대응하도록 설계되어 있다.

\subsection{엑셀 또는 스프레드시트에서 복사하여 붙여넣기}

\begin{center}
\includegraphics[scale=.35]{screenshot002}
\end{center}

\tblrdoumi 를 실행한 후에 다음 조치를 먼저 취한다.
\begin{itemize}
\item \hcrkey{ESC}를 눌러서 \guielem{테이블 영역}으로 들어간 후에 \hcrkey{F9}를 누르면 괘선과 행\.열\.셀 설정이 싹 비워진다. (\hcrkey{F9}는 포커스가 테이블 영역에 있어야 효과가 있다.)
\item {} [기본 설정] 탭의 Table Size: 부분에서 \uline{행과 열을 복사해올 행과 열의 수에 맞추어 두어야} 한다. 이를테면 6행 3열을 만들고 보라색 \fbox{적용} 버튼을 누른다.
\end{itemize}

\noindent 
스프레드시트에서 전체 영역을 선택하여 복사(\hcrkeys[+]{Ctrl, C})한 다음, \tblrdoumi 의 (1,1) 셀에서 붙여넣기(\hcrkeys[+]{Ctrl, V})한다.

\subsection{엑셀 파일 들여오기}

상단의 초록색 버튼은 엑셀 파일 들여오기이고 그 옆의 파란 버튼을 눌러 메뉴를 열어도 [엑셀 들여오기]가 있다.

이 기능은 말 그대로 엑셀 스프레드시트(\path{*.xlsx})를 \tblrdoumi\ 표로 불러들이는 것이다. 우리는 지금 표의 본체를 가져오는 데 관심이 있지만 실제 이 기능은 병합 셀을 해석한다든가 하는 등 몇 가지 알찬 기능을 더 가지고 있다. 함수가 포함된 셀은 물론 계산 결과를 가져온다.

다만 기억해야 할 것은 Excel의 모든 정보를 다 가져올 수는 없다는 것이다. 예컨대 엑셀에서 그은 괘선 정보는 전혀 가져오지 않는다.

이 기능에 의해 가져온 데이터는 텍스트는 중괄호로 감싸고 숫자는 그대로 노출시킨 형태이다. 숫자인가 아닌가는 Excel의 설정에 따른다.

\section{파일로 저장하기, 불러오기}

파란색 메뉴를 열면 [저장하기]와 [불러오기]가 있다. 각각 \hcrkeys[+]{Ctrl, S}와 \hcrkeys[+]{Ctrl, O} 단축키로 실행할 수 있다.

이 기능은 주로 작업의 연속성을 위하여 마련한 것이었는데 \ksemph{현재의 작업 상태를 그대로 저장}했다가
불러와서 이전 작업을 이어가자는 것이었다.

파일은 JSON 형식의 독자 포맷 \path{*.tblr}이다. Windows에서는 설치 프로그램을 이용하여 설치한 경우에 더블클릭하여 \tblrdoumi 에서 열 수 있다. 즉, \ksparensmall{\LaTeX\ 코드가 들어 있기는 하지만} 즉시 문서에 넣을 수 있는 \LaTeX\ 출력이 아니다.

\tblrdoumi 는 마지막에 접근한 폴더를 기억하게 하여 불러올 파일을 찾기 위해 이리저리 폴더를 뒤적여야 하는 수고를 줄이려고 했다.


\chapter{셀(Cells)}

\section{셀의 병합과 병합 풀기}

\env{tabular}에서 \cmdprint{\multicolumn}과 \pkg{multirow}에 의한 \cmdprint{\multirow}로 구현하던 것이다. 옛날 코드를 감상하고 가자.

\begin{widetext}
\setrighthandwidth{.4\linewidth}
\begin{exampleside}
\begin{tabular}{l|c|r}
\textbf{Value 1} & \textbf{Value 2} & \textbf{Value 3}\\ \cline{1-2}
\multicolumn{2}{c|}{$\alpha\beta$} & $\gamma$ \\ \hline
\multicolumn{2}{c|}{\multirow{2}{*}{1234}} & a\\
\multicolumn{2}{c|}{} & b\\ \hline
\multirow{2}*{3} & 23.113231 & c\\
 & 25.113231 & d\\
\end{tabular}
\end{exampleside}
\end{widetext}

표를 그리다가 곧잘 \verb|&|의 위치를 잃어버리거나 개수를 맞추지 못하거나 하는 일이 일어나는 것은 저 번잡스러운 코드 때문이었다. 표의 본체에 추가적인 명령어를 계속 입력해 넣어서 원하는 결과를 얻는 방식으로는 어쩔 도리가 없었다고 생각한다.

\tblrdoumi 는 셀의 병합을 매우 우아하게 처리한다. 즉 둘 이상의 인접한 셀을 선택한 뒤에 우클릭하면 아래 그림과 같이 \texttt{선택 셀 병합} 항목이 활성화된다. 단축키 \hcrkeys[+]{Ctrl, M}(\uline{M}erge)로 같은 일을 할 수 있다.

\begin{center}
\includegraphics[scale=.35]{screenshot003}
\end{center}

앞에서 보인 예제처럼 병합하려면,
\begin{enumerate}[\quad1.]\oblivoirlist
\item (2,1) 셀과 (2,2) 셀을 선택하여 병합한다.
\item (3,1), (3,2), (4,1), (4,2)를 한꺼번에 선택하여 병합한다.
\item (5,1)과 (6,1)을 병합한다.
\end{enumerate}
와 같은 절차를 따르면 될 것이다.

\guielem{테이블 영역}에 병합된 모양이 잘 나타나므로 그 뒤로는 신경쓸 것이 없다. \LaTeX\ 코드, 즉 \env{tblr} 코드는 \tblrdoumi 가 알아서 잘 생성해준다.

병합 풀기는 병합된 셀에서만 (우클릭 메뉴에서) 활성화되며, \hcrkeys[+]{Ctrl, U}(\uline{U}nmerge) 단축키로 동작한다. 병합되지 않은 셀에서는 아무 일도 일어나지 않는다.

\begin{Chamgo}
내용물이 있는 두 셀을 병합했을 때 왼쪽 위에 있는 셀(host 셀)의 내용만이 \LaTeX\ 코드에 반영된다. 그렇지만 병합된 셀의 내용도 앱이 기억하고 있으므로 병합을 풀면 그 자리에 나타날 것이다.

내부적으로 병합의 처리는 \cmdprint{\SetCell}에 의하지 않고 설정부의 \texttt{cell} 설정에 의한다. 앱의 \guielem{설정 영역}의 Cell: 파트에 직접 병합 명령을 써넣거나 표의 본체 셀에 \cmdprint{\SetCell} 명령을 주면 \LaTeX\ 출력은 나타나겠지만 앱의 ``병합 정보''로는 다루어지지 않는다. 병합 정보는 별도로 관리하고 있으므로 병합에 대해서는 앱의 지침을 따르는 것이 여러 가지로 좋다.
\end{Chamgo}

\cmdprint{\multicolumn}이나 \cmdprint{\multirow}와는 비교할 수 없을 정도로 일관성있고 알아보고 조작하기 쉽게 되었다고 본다. 이러한 조작을 행하더라도 표의 본체는 조금도 변화가 없다.

\section{셀 서식}

\subsection{대화창을 통한 기본 서식}

셀의 기본 서식을 설정하는 대화창은 \hcrkey{F5}이다. 셀에서 우클릭하여 ``셀 설정''이라는 항목을 선택해도 같다.

\begin{center}
\includegraphics[scale=.5]{screenshot004}
\end{center}

\noindent 이 대화창을 통하여 설정할 수 있는 것은 다음 항목이다.
\begin{description}
\item [배경색] \texttt{bg=<color>}.
\item [글자색] \texttt{fg=<color>}.
\item [halign] \texttt{r}, \texttt{c}, \texttt{l}, \texttt{j}.
\item [valign] \texttt{t}, \texttt{m}, \texttt{b}, \texttt{h}, \texttt{f}.
\item [mode] \texttt{text}, \texttt{math}, \texttt{imath}, \texttt{dmath}.
\end{description}

팔레트 아이콘을 누르면 \pkg{ninecolors} 색상 선택 대화창이 열린다. 여기서 색상을 선택할 때 그 오른쪽의 콤보박스가 배경색(default)으로 되어 있다면 
\begin{verbatim}
  bg=<color>
\end{verbatim}
가 생성된다.
\verb|bg=| 부분을 생략하고 싶다면 [직접 입력]을 골라두고 색상을 선택한다.

[가로 유지]라고 된 것은 \verb|halign|에 아무 것도 선택하지 않는다는 의미이다. 콤보박스를 열어서 가로 정렬자를 고를 수 있다.

[세로 유지]라고 된 것은 \verb|valign|에 아무 것도 선택하지 않는다는 의미이다. 콤보박스를 열어서 세로 정렬자를 고를 수 있다.

마지막의 [text]는 \verb|mode|를 선택하게 하는 것이다. \verb|math|, \verb|imath| (inline math), \verb|dmath| (display math) 중에서 고른다.

이 대화창에서 OK를 누르면 선택 사항을 반영하여, 오른쪽 \guielem{설정 패널}의 \guielem{Cell Format:} 편집창에 다음과 비슷한 텍스트가 입력될 것이다.
\begin{verbatim}
  cell{2}{3} = {c, h, mode=math, bg=purple2, fg=yellow9}
\end{verbatim}
아무 것도 선택한 것이 없다면 빈 중괄호로 들어온다. 그리고 halign 선택사항을 반영해서 \guielem{테이블 영역}의 해당 셀의 텍스트 정렬 상태가 시각적으로 바뀐다(단, \verb|c|, \verb|l|, \verb|r|만 반영함).

맨 왼쪽의 작은 에디트 영역에 원한다면 추가 설정, 색상명 등을 직접 써넣을 수 있다. 그러나 이곳은 너무 좁으니까 다음 소절에서 말하는 대로 설정 패널의 에디터에서 하는 것이 좋을 것이다.

\subsection{설정 패널에서 직접 수정}

색상을 꼭 \pkg{ninecolors}만 쓰라는 법은 없을 것이다. 그리고 \opt{font=}, \opt{cmd=}과 같은 추가적인 설정 옵션이 더 필요할 때가 있다. 이것은 Cell Format: 설정 행을 직접 수정하여야 한다. 

\subsection{설정 패널의 편집 영역}

\guielem{설정 패널}이라고 부르는 영역은 가리거나 불러낼 수 있다. 화면 오른쪽 상단의 검은색 버튼을 이용하거나 \hcrkeys[+]{Ctrl,=}) 토글 단축키를 이용한다.

이 부분에 상단의 몇몇 설정을 제외하고 네 개의 \guielem{설정 편집 영역}이 있다. \texttt{Lines}, \texttt{Row/Col Format}, \texttt{Cell Format}, \texttt{Generic}이 각각 있는데, 차지하는 공간이 좁아서 몇 줄 보이지 않을 것이다.
이 네 개의 작은 편집기는 떼어내거나 떠다니게 하거나 겹쳐서 탭 배열로 만들 수 있다. 만약 내용이 많아서 한눈에 보이지 않으면 떼어낸 다음 창을 키워서 불 수 있다. 다음 그림은 네 개의 편집 영역을 탭으로 겹쳐쌓은 경우이다.

\begin{center}
\includegraphics[scale=.35]{screenshot005}
\end{center}

\subsection{고급 셀 설정}

앞서 \hcrkey{F5} 셀 설정 대화창에서 [OK]를 누르면 어떤 내용이 Cell Format 편집기에
등록된다고 하였다.
이 셀은 ``설정문이 있는 셀''이라고 하여 셀 오른쪽 위에 조그맣고 파란 표지가 하나 붙는다.
그리고 포커스를 이리저리 움직이다가 이 셀에 오게 되면 해당 설정문의 색이 파랗게 변한다. 위의
스크린샷에서 \texttt{cell\{1\}\{2\}=...} 부분의 색이 변한 것은 지금 포커스가 (1,2)에 있으며
그 셀에 대한 설정문이 이 자리에 존재한다는 의미이다.

셀 선택 후에 우클릭하면 ``셀 설정''과 별개로 ``셀 고급 설정''이라는 항목이 있다. 이 항목은 \hcrkeys[+]{Shift,F5}로 호출할 수 있는데, 하는 일은 다음과 같다.
\begin{itemize} \oblivoirlist
\item 이미 설정이 부여된(별표가 붙은) 셀에서 호출하면 해당 설정문으로 포커스가 이동하여 수정할 수 있게 한다.
\item 설정이 부여되지 않은 셀에서 호출하면 해당 셀의 좌표를 이용하여 Cell Format 편집 영역에 새로운 문장을 만들고 편집할 수 있게 한다.
\end{itemize}
그러므로 아까 \hcrkey{F5}로 만든 셀 설정 선언문이 있다면 같은 셀에서 셀 고급 설정을 눌러서 해당 설정문으로 간 다음, 추가할 내용을 적어넣으면 되는 것이다.

\guielem{테이블 영역}과 \guielem{설정 편집 영역} 사이를 오가면서 설정을 편집할 수 있게 되었다고 본다.

\begin{Chamgo}
설정 편집기들이 앱으로부터 떨어져나와 떠다니게 되면 이 편집기 안에서 \hcrkey{ESC}를 눌러도 \guielem{테이블 영역}으로 초점이 이동하지 않는다. 그러나 편집기 안에 붙어 있는 상태라면 \hcrkey{ESC}를 통해서 포커스를 이동할 수 있다.
\end{Chamgo}

\begin{Chamgo}
설정 편집기에서 각 지시문은 행 단위로 입력한다. 행 끝에 쉼표(\verb|,|)를 치지 않아도 앱 자신이 \LaTeX\ 코드를 생성하면서 필요한 쉼표를 추가한다. 이렇게 설계한 이유는 \env{tblr} 설정부에서 콤마를 잘못 치거나 빠뜨려서 곤란한 경험을 한 적이 있어서였다.

따라서 하나의 설정문을 두 줄 이상으로 나누어서 입력하지 않아야 한다. 입력된 행의 마지막에 자동으로 쉼표가 붙기 때문이다.
\end{Chamgo}

\chapter{행(rows)과 열(columns)}

행과 열 설정은 셀 설정과 본질적으로 같다. ``설정이 부여된 셀''에 조그마한 표지가 나타났듯이
``설정이 부여된 행/열''에는 표의 헤더 영역(행번호 또는 열번호가 디스플레이되는 영역)의 행/열 번호 옆에 조그마한 동그라미가 붙는다.

단축키를 기억해두면 편리할 것이다. 물론 우클릭 콘텍스트 메뉴를 통해 접근할 수 있다.

\begin{description}
\item [{\hcrkey{F6}}] 행 설정 (다이얼로그).
\item [{\hcrkeys[+]{Ctrl,F6}}] 열 설정 (다이얼로그).
\item [{\hcrkeys[+]{Shift,F6}}] 행 고급 설정.
\item [{\hcrkeys[+]{Shift,Ctrl,F6}}] 열 고급 설정.
\end{description}

행과 열에 대해서도 배경색\.글자색을 지정한다. \pkg{ninecolors} 색상 선택 패널을 이용할 수 있다.

\begin{center}
\includegraphics[scale=.5]{screenshot010}
\end{center}

``고급 설정''을 통해서 설정 편집 영역의 에디터로 이동하는 것은 동일하다.

\begin{Chamgo}
수작업으로 여러 행에 대하여 예를 들면 \texttt{row\{1-Z\}}와 같은 표현을 쓴 경우에 
해당 설정문이 해당 행에 적용되어 있음을 알려주기는 하지만 예컨대 마지막 행에 대하여 
``행 설정''을 선택한다고 해서 \texttt{Z}로 나타나지는 않는다.
\end{Chamgo}

\paragraph{Generic 편집 영역}

설정 에디터 영역의 하나인 Generic은, 괘선 설정, 행\.열\.셀 설정에 해당하지 않지만 \env{tblr}의 설정부에 오는
모든 내용을 적어넣을 수 있다. \opt{manual} 모드의 \verb|width=| 설정도 이 영역에 적는다.

그밖에 \verb|stretch=|, \verb|rowspec=| 등의 문장을 활용하는 경우에도 이 영역을 이용할 수 있다.

\tblrdoumi 는 \opt{colspec}을 매우 중요하게 다룬다. 그러나 \opt{rowspec}에는 거의 관심이 없다. 그러므로 굳이 \opt{rowspec}으로 표 형태를 작성하려 한다면 Generic 부분에 해당 선언을 적어넣도록 하자.

\chapter{Colspec}

\section{colspec을 다루는 방식}

아래 그림은 \tblrdoumi 가 colspec을 설정하는 데 관련된 부분을 표시한 것이다.
\begin{center}
\includegraphics[scale=.4]{screenshot006}
\end{center}

최상단의 \guielem{colspec 미리보기} 영역은 편집불가능하고 현재 활성화된
colspec을 보여준다. 이것은 그대로 하단 아래의 \guielem{LaTeX 미리보기}에
나타난다.

중요한 것은 \guielem{설정 패널}의 세 개의 라디오버튼이다. 앱이 내부적으로 ``colspec mode''라고 부르는 것을 선택할 수 있다.

\section{colspec mode}

\subsection{colspec mode의 종류}

각각 다음과 같은 특징을 갖는다.
\begin{description}
\item [natural] 열의 수만큼 \verb|l|을 넣는다. 모든 셀이 왼쪽 정렬 상태로 나타난다.
\item [linewidth] 표의 width를 \cmdprint{\linewidth}로 해놓고 열의 수만큼 \verb|X[c,m]|을 넣는다. 모든 셀이 가운데 정렬 상태로 나타난다.
\item [manual] 열의 수를 앱이 조절하지 않으며 사용자가 직접 colspec을 편집하거나 수정할 수 있다. 다만 사용자가 지정한 열의 수가 표의 열 수와 맞지 않으면 경고를 보이며, \verb|X| 열지정자가 사용되었는데 \texttt{width=} 행이 없으면 직접 써넣을지를 묻는다. 이때는 Generic: 부분에 설정행이 생기므로 이를 수정하여 표 전체 폭을 조절할 수 있다.
\end{description}

\opt{natural}과 \opt{linewidth} 모드는 ``완성되지 않은 표 입력 작업''을 진행하는 과정에서 열의 수 문제로 생기는 불일치를 막기 위해 설정한 것이다. 이때에는 새로운 열을 추가하거나 삭제하는 일도 많아서 매번 colspec의 열지정자 수를 세어야 할 때가 있기 때문에 이를 피하려고 하는 것이다.

그러므로 최종적인 표의 모양은 \opt{manual} 모드에서 편집하는 colspec이 되는 것이다.

\subsection{manual 모드}

\opt{manual} 모드에서 중요한 것은 ``현재 \texttt{manual colspec}''으로 지정된
문자열이 존재하는가이다. 이 문자열은 라디오버튼 바로 아래의 manual colspec:으로 되어 있는
부분에 적혀 있는데 여기어 어떤 값이 있든 상관없이 \opt{natural}, \opt{linewidth}에서는
나름의 규칙으로 열 수만큼 colspec을 자체 생성한다. 이곳은 비어 있을 수도 있고 (이전 어떤 작업의 결과) 비어 있지 않을 수도 있다.

\begin{itemize} \oblivoirlist
\item \guielem{라디오버튼}에서 \opt{manual}을 선택한다. 그러면 앞서 말한 manual colspec의 문자열이 실제 colspec으로 활성화된다(상단의 미리보기를 통해 확인 가능).
\item 만약 manual colspec에 있던 문자열을 실제 colspec으로 가져올 때 열의 수가 맞지 않거나 이해할 수 없는 지정자가 들어 있으면 경고를 보인다. ``이후로 무시하라''고 지정할 수도 있다.
\end{itemize}

만약 \opt{manual} 모드에 놓여 있다면 오른쪽의 [colspec 편집] 버튼을 눌러서 현재 colspec 문자열을 수정하거나 편집할 수 있다. 앱의 불일치 체크에서 유효한 것으로 간주하는 열 지정자는 \opt{r}, \opt{c}, \opt{l}, \opt{X}, \opt{Q}이다. \opt{X}는 대괄호로 인자를 더 취할 수 있는데 이것이 쓰이면 width 설정이 있는지만 검사한다. \opt{Q}는 대괄호 인자를 취할 수 있는데 이 안에 오는 내용에 대해서는 에러 체크를 하지 않는다.

\pkg{tabularray}에서 제공하는 기타 열형식(예: \verb|S|)은 유효하지 않다고 경고할 것이지만 무시하게 하고 자기 책임으로 진행하면 된다. 새로운 칼럼 타입을 정의한 경우에도 그러하다. 
그리고 이른바 ``별표 반복문''은 반복 횟수를 계산하여 열 개수와 비교하므로 \opt{manual} 모드에서는 그대로 쓸 수 있다.

\paragraph{F10 단축키}
포커스가 \guielem{테이블 영역}에 있을 때 \hcrkey{F10}을 누르면 \opt{natural}이나 \opt{linewidth} 모드에서 만들어진 colspec을 \opt{manual colspec}으로 보내고 모드를 \opt{manual}로 바꾼다.

예를 들어 \opt{linewidth} 모드로 시작하여 6열 표를 완성하였다면 현재 colspec 미리보기에는 다음 문자열이 있을 것이다.
\begin{verbatim}
X[c,m] X[c,m] X[c,m] X[c,m] X[c,m] X[c,m]
\end{verbatim}
\hcrkey{F10}을 누른 결과는 모드가 \opt{manual}로 바뀌고 manual colspec이 이것으로 교체되는 것이다.
여기서 옆의 [colspec 편집] 버튼을 누르거나 \hcrkey{F12}를 누르면 이 colspec 문자열을 수정\.편집할 수 있다.

\paragraph{colspec 편집 버튼}
\hcrkey{F12} 키는 편집 버튼을 누른 것과 같다. 이 버튼을 누르면 일어나는 일은 현재 어떤 colspec 모드에 있느냐에 따라 다르다.
\begin{itemize}
\item \opt{natural} 모드에 있을 때: 현재의 \verb|l|로 이루어진 colspec 열을 그대로 가져와서 manual colspec으로 만든 다음 colspec 편집창을 열어준다. 여기서 수정한 결과가 manual colspec이 되고 모드는 \opt{manual}이 된다.
\item \opt{linewidth} 모드에 있을 때: 현재의 \verb|X[c,m]|으로 이루어진 colspec 열을 가져오되, 여기에 (현재 테이블 영역의 시각적인) 열 너비의 상대 길이를 계산하여 \verb|X|의 (\verb|wd|에 해당하는) 인자로 넣어준다. 즉 마우스로 조절한 열의 상대 길이를 반영하여 manual colspec을 만든 다음 colspec 편집창을 열어준다. 여기서 수정한 결과를 manual colspec으로 하고 모드는 \opt{manual}이 된다.
\item \opt{manual} 모드에 있을 때: 현재의 manual colspec 문자열로써 colspec 편집창을 열어준다.
\end{itemize}

\begin{Chamgo}
\opt{manual} 모드 자체에는 \texttt{width=\bs linewidth}가 설정되지 않으므로 
이것을 자동으로 써넣을지 묻는 대화창이 열릴 수 있다.
\end{Chamgo}

\subsection{열의 상대 너비}

\guielem{테이블 영역}에서 열 헤더의 가장자리를 마우스로 끌어서 열의 (시각적) 폭을 조절할 수 있다. 

\begin{center}
\includegraphics[scale=.4]{screenshot007}
\end{center}

둘 이상의 열을 선택하여 마우스로 헤더를 조작하면 (원래 폭에 비례해서 변하는 게 아니라) 같은 폭으로 변화한다.

우클릭하여 뜨는 메뉴 항목 중의 ``열 너비 초기화''를 클릭하면 모든 열이 처음 상태로 돌아간다.

이 상태에서도 \opt{natural} 모드나 \opt{linewidth} 모드에서는 열의 폭에 관계없이 colspec이 동작한다. 그런데 이 상대 너비를 colspec에 반영할 방법이 있다.
상대 너비는 (시각적으로) 가장 좁은 폭의 컬럼을 $1$로 하고 다른 열의 현재 길이를 이것의 배수로 환산하는 것이다.

\paragraph{콘텍스트 메뉴에서}
``열 상대 너비 반영''을 클릭하면 (현재 모드에 상관없이) 열 상대 너비를 계산하여 \verb|X| 열로 만들고 \opt{manual} 모드로 전환한다. 위 예시 그림의 경우는 다음과 같은 새로운 manual colspec 문자열이 생성된다.
\begin{verbatim}
   X[1,c,m] X[2.2,c,m] X[2.7,c,m]
\end{verbatim}

\paragraph{colspec 편집 버튼}
버튼 또는 \hcrkey{F12} 동작은 현재 모드가 \opt{linewidth}일 때만 열의 상대 너비를 계산하여 \opt{manual}로 가져온다.

\paragraph{manual 모드에서 열 너비의 변경}
\guielem{테이블 영역}에서 마우스를 끌어서 열 너비를 재조정하였다면, 이를 반영하기 위해서는 우클릭 후 ``열 상대 너비 반영''을 클릭하여야 한다.

\begin{Chamgo}
\opt{manual} colspec 편집에서 수작업으로 숫자를 변경하는 것은 최종 \LaTeX\ 코드 생성에 영향을 준다. 최종 출력으로 삼는 것은 수작업 수정의 결과가 되겠지만, 수작업 수정이 시각적으로 테이블 영역 열의 폭에 반영되는 것은 아니다.
\end{Chamgo}

\chapter{괘선}

\section{초기 상태의 괘선 코드}

맨처음 앱을 실행하면 \guielem{Lines} 에디터에 다음과 같은 내용이 들어가 있다.
\begin{verbatim}
   hline{1,Z} = {0.8pt, solid}
   hline{2} = {solid}
   vline{2-Y} = {solid}
\end{verbatim}
그 의미는 명확하다. 가로 괘선에서 첫 줄과 마지막 줄은 굵게, 둘째 줄은 보통 가늘기 선을 하나 그리고, 
세로 괘선은 처음과 마지막을 제외한 모든 선을 기본 굵기로 그리라는 것이다.
이 형식은 \guielem{테이블 영역}에 시각적으로 나타난다.

이 샘플은 \tblrdoumi\ 앱의 기능 소개. 괘선 설정 영역을 어떻게 이용할 수 있는지 보이기 위한 것이다. 그러므로 자신의 표를 작성할 때는 일단 지우고 시작하는 것이 좋다. 앞서 \hcrkey{F9} 단축키를 소개했지만 이 에디터 영역에 커서를 두고 \hcrkeys[+]{Ctrl, A}, \hcrkey{Del}를 차례로 누르면 다 지워지겠지.

테이블 영역에 나타나는 시각적인 괘선 긋기 효과는 단순히 작업을 보조하기 위하여 ``여기에 괘선이 있다''는 정도를 알아보기 위한 목적이다. WysiWyg 프로그램처럼 괘선의 모양이 출력 형태를 보여주기 위한 것이 아니라는 사실을 강조해둔다.

\section{괘선 설정 다이얼로그}

\subsection{선택한 범위의 괘선}

\guielem{테이블 영역}에서 괘선을 설정하는 기본적인 순서는 다음과 같다.
\begin{enumerate}[\quad 1.] \oblivoirlist
\item 괘선의 범위를 선택한다.
\begin{enumerate}[\quad 1)] \oblivoirlist
\item 셀 선택으로 괘선의 범위를 정하는데 가로 괘선은 선택된 셀의 상단, 세로 괘선은 선택된 셀의 왼쪽 괘선을 의미한다.
\item 마지막 가로 괘선과 마지막 세로 괘선\ksparensmall{외곽선}은 \verb|Z| 괘선이라 하여 별도로 선택하여 설정할 수 있다.
\item 행\.열 전체가 선택되면 \cmdprint{\hline}과 같이 전 범위에 걸치는 가로\.세로 괘선을 가리킨다. 반면 일부 셀만 선택하면 \cmdprint{\cline}처럼 일부 셀에만 해당하는 범위에 걸친다.
\item 마지막 괘선에 대해서도 같다. 
\end{enumerate}
\item 우클릭 또는 단축키를 이용하여 괘선 설정 대화창을 연다.
\begin{itemize} \oblivoirlist
\item \hcrkeys[+]{Alt, H}(\uline{h}orizontal) 가로 괘선
\item \hcrkeys[+]{Alt, Z} 마지막 가로 괘선
\item \hcrkeys[+]{Alt, V}(\uline{v}ertical) 세로 괘선
\item \hcrkeys[+]{Shift, Alt, Z} 또는 \hcrkeys[+]{Ctrl, Alt, Z} 마지막 세로 괘선
\item 참고로, Mac 앱에서는 \hcrkey{Alt} 자리에 Option({\applesymbolfont{1.5}⌥})이 아니라 Control({\applesymbolfont{1.3}⌃}) 키가 배당되어 있다.
\end{itemize}
\item 괘선의 색상, 굵기, 점선 여부, 이중선 여부를 선택한다.
\end{enumerate}

\bigskip

괘선 설정 대화창은 다음처럼 생겼다.
\begin{center}
\includegraphics[scale=.5]{screenshot009}
\end{center}

맨 왼쪽 에디트 영역에 원하는 것을 써넣을 수 있다. 그 다음 팔레트 아이콘은 \pkg{ninecolors} 색상 패널을 부르는 것이다. 색상을 고를 수 있다. 행\.열\.셀 설정 대화창에 붙어 있는 색상 고르기 패널과 동일하다.

콤보박스는 \opt{solid}, \opt{dashed}, \opt{dotted} 중에서 고르는 것이다.

오른쪽에 두 개의 체크박스가 있는데 각각 \opt{double}과 \opt{thick}이다. 엄밀히 말해서 \pkg{tabularray}는 \verb|double|이나 \verb|thick| 같은 키워드를 바로 인식하지 않는다.
그래서 이곳을 체크하면 상황에 맞게 적절한 구문으로 \guielem{Lines}에 지시문을 넣어준다.

이중선\ksparensmall{남용은 억제하는 편이 좋다고 생각하지만} 그리기가 꼭 필요할 때 의외로 편리하다. double line은 파란 줄 하나로 표시된다.

굵기는 다양하게 제공하지 못하고 단지 \opt{thick}을 체크하면 \opt{0.8pt}인데 이 값은 에디터 영역에서 원하는 대로 수정할 수 있다.

dashed, dotted는 비슷한 모양으로 테이블 영역에서 보이도록 애썼다.

\subsection{마지막 괘선}

마지막 괘선은 위에서 소개한 단축키로도 설정할 수 있지만 다른 괘선과 \ksemph{함께} 설정할 수도 있다. 
콘텍스트 메뉴의 ``가로괘선+Z''라는 것과 ``세로괘선+Z''라는 것이 이에 해당하고 각각 \hcrkeys[+]{Alt,Shift,H}와 \hcrkeys[+]{Alt,Shift,V} 단축키가 할당되어 있다. 
현재 선택된 범위의 셀 영역에서 상단 또는 왼쪽 괘선과 최후 괘선을 동일하게 설정하는 것이다.

\section{실제 예}\label{sec:exam}

앞서 보았던 것과 비슷한 다음 도표의 괘선을 실제로 그어보자.

\begin{center}
\begin{tabular}{l|c|r}
\textbf{Value 1} & \textbf{Value 2} & \textbf{Value 3}\\ \cline{1-2}
\multicolumn{2}{c|}{$\alpha\beta$} & $\gamma$ \\ \cline{2-3}
\multicolumn{2}{c|}{\multirow{2}{*}{1234}} & a\\
\multicolumn{2}{c|}{} & b\\ \hline
\multirow{2}*{3} & 23.113231 & c\\
 & 25.113231 & d\\ \hline
\end{tabular}
\end{center}

\begin{Chamgo}
앞서 셀 병합(merge)에 대해 말했지만 괘선을 그을 때 어떤 괘선에 대해서는 병합된 셀이 정확한 위치를 알려주지 않을 수 있으므로 모든 병합을 풀고 괘선 작업을 하는 것이 좋다. 혹은, 괘선 작업을 병합 작업보다 먼저 하는 것이 좋다.
\end{Chamgo}

우선 복잡한 병합 셀을 모두 푼 다음에 괘선 작업을 하기로 하자. \uline{(2,1)}--(2,2) 병합 셀과, \uline{(3,1)}--(4,2) 병합셀, 그리고 \uline{(5,1)}--(6,1) 병합셀에서 \hcrkeys[+]{Ctrl,U}를 눌러 병합을 모두 해제한다.

\paragraph{세로 괘선} 2, 3번 열의 왼쪽 선이다.

2, 3번열을 선택한 다음 \hcrkeys[+]{Alt, V}를 눌러서 다이얼로그를 띄운 다음 그냥 [OK]를 눌러 닫는다. 그러면 \guielem{Lines} 에디터 영역에
\begin{verbatim}
  vline{2,3} = {}
\end{verbatim}
이라는 행이 생기고 세로 선이 \guielem{테이블 영역}에 나타날 것이다. 저 빈 중괄호는 실은 \verb|{solid}|와 동일한 의미이기 때문에 이렇게 된다.

\paragraph{가로 괘선} 우선 통짜로 전체를 그리는 가로 괘선은 5번 행의 상단과 마지막 괘선(Z)이다.
그러므로 5번 행을 전체 선택한 후 ``가로 괘선 + Z''를 콘텍스트 메뉴에서 선택(단축키로는 \hcrkeys[+]{Shift, Alt, H})한다.
색상도 굵기도 선택할 필요가 없으므로 그냥 [OK]를 눌러서 빈 중괄호로 만들어 넣어도 된다.

그 다음은 일부 셀 범위에만 걸치는 가로 괘선이다.
먼저 (2,1)--(2,2) 범위를 셀 선택한 후에 콘텍스트의 ``가로 괘선''(단축키로는 \hcrkeys[+]{Alt, H})을 선택하여 (앞서 설명한 바와 같이) [OK]를 누른다.

다음으로 (3,2)--(3,3) 범위를 셀 선택한 후에 같은 작업을 반복한다.

이것으로 예제 표의 모든 괘선이 잘 그려졌다. 최종적으로 \guielem{Lines} 에디터 창에 들어가 있을 텍스트는 다음과 같을 것이다.

\begin{center}
\includegraphics[scale=.4]{screenshot011}
\end{center}

\begin{verbatim}
   vline{2-3} = {}
   hline{5,Z} = {}
   hline{2} = {1-2}{}
   hline{3} = {2-3}{}
\end{verbatim}

\section{직접 수정}

괘선에 있어서도 복잡한 경우에는 이런저런 설정을 추가해야 할 수 있다. 
또는 앞서 설명한 GUI 방식의 설정이 귀찮고 단번에 그냥 텍스트로 입력하고 싶을 때가 있다.
굵은 괘선의 두께를 조절하거나, 색상을 \pkg{ninecolors}가 아닌 다른 것으로 설정하려 하거나.

이럴 때는 설정 편집 영역의 \guielem{Lines} 에디터를 이용하면 되겠다.

\chapter{Long과 Tall}

\tblrdoumi 는 기본 \env{tblr} 환경과 함께 \env{longtblr}, \env{talltblr} 환경을 지원한다. 

\guielem{설정 패널}의 첫 항목 \guielem{테이블 환경} 콤보박스에서 선택한다. 그러면 [long/tall 설정]이라는 탭이 활성화된다.

\opt{long}일 때는 ``선택 옵션''과 ``추가 설정''이라는 두 입력 에디터가 들어 있는데 각각 어떤 것을 적어넣어야 하는지 간략한 보기를 들고 있다.

\opt{tall}일 때는 ``선택 옵션'' 한 종류의 입력 에디터밖에 없는 대신 \verb|note{}={}|라는 문장을 예시하고 있다. 셀 텍스트 입력시에 \cmdprint{\TblrNote} 명령은 직접 적어넣어야 한다. 

특히 엑셀 데이터를 불러온다든가 하면 많은 경우에 한 페이지에 다 들어가지 않는 긴 표일 것이다.
\env{longtblr} 지원은 이런 상황에 편리하게 대응할 수 있게 한 것이다.
그러나 표 환경 밖에서 해야 하는 \env{longtblr} 설정을 \tblrdoumi 가 모두 관여하지는 않는다.

\chapter{출력}

표 작성 작업이 끝났다면 이 결과를 어떻게 활용할 것인가? 
\verb|.tblr| 포맷의 파일은 그대로 \LaTeX\ 문서에 삽입하거나 할 수 있는 것이 아니므로 생성한 \LaTeX\ 코드를 이용할 방법이 있어야 할 것인데,
원래 이 프로그램을 처음 설계할 당시에는 단순히 ``클립보드로 복사''라는 한 가지 출력 방법밖에 고려하지 않았다.

\section{클립보드로 복사}

가장 간단하지만 효과적인 방법이다. \hcrkey{F4} 단축키 또는 파란색 메뉴에서 ``클립보드로 복사''를 선택하면 왼쪽 아래에 잠깐 ``복사되었음''이라는 메시지가 나타난다. 

그리고 작성중인 \LaTeX\ 문서에서 붙여넣기하면 완성한 코드가 그대로 삽입되는 것이다.

\section{\texttt{tex} 파일로 내보내기}

파란색 메뉴를 열면 ``tex 내보내기 [표만]''이라는 항목이 있다. 단축키는 \hcrkeys[+]{Ctrl,T}(\uline{T}eX)이다.
말 그대로 현재 앱이 완성한 \LaTeX\ 코드를 \path|.tex| 파일로 저장하는 것이다. 작성 중인 \LaTeX\ 문서에서는 \cmdprint{\input} 명령을 써서 포함할 수 있다.

클립보드로 복사와 \opt{tex} 파일로 내보내기는 본질적으로 같은 것이다. 파일로 저장하느냐 않느냐의 차이이다.

\pageref{sec:exam}페이지에서 작성했던 표로써 생성되는 코드는 다음과 같다.

\begin{exampleonly}
\begin{tblr}{
  colspec = {l c r},
  vline{2-3} = {},
  hline{5,Z} = {},
  hline{2} = {1-2}{},
  hline{3} = {2-3}{},
  row{1} = {font=\bfseries},
  cell{2}{1} = {r=1, c=2}{},
  cell{3}{1} = {r=2, c=2}{},
  cell{5}{1} = {r=2, c=1}{},
}
  Value 1 & Value 2 & Value 3 \\
  $\alpha\beta$ &  & $gamma$ \\
  1234 &  & a \\
   &  & b \\
  3 & 23.113231 & c \\
   & 25.113231 & d \\
\end{tblr}
\end{exampleonly}

\section{클래스를 지정하여 tex 파일로 내보내기}

표만을 \verb|.tex|으로 내보내는 것이 아니라 컴파일 가능한 형식의 \texttt{tex} 문서로 내보내는 것이다. 
파란색 메뉴의 ``tex 내보내기''에서 서브 메뉴를 열어보면 선택할 수 있는 클래스로 \texttt{oblivoir}, \texttt{standalone}, \texttt{article}이 있고 
그 가운데 \texttt{oblivoir} 문서로 내보내기의 단축키가 \hcrkeys[+]{Shift,Ctrl,T}인 것을
알 수 있다.

그 아래의 \opt{class selection...} 항목은 특정 클래스를 직접 지정할 수 있다.

이 기능은 \verb|\documentclass|와 \verb|\begin{document}|, \verb|\end{document}|를 모두 갖춘 컴파일 가능한 파일을 생성한다. 프리앰블은 다음과 같이 되어 있다.

\begin{verbatim}
\usepackage{tabularray}
\usepackage{xcolor,ninecolors}
\IfFileExists{mytblrpreamble.tex}{\input{mytblrpreamble.tex}}{}
\end{verbatim}

\noindent 이 코드가 의미하는 것은 다양한 패키지나 라이브러리를 활용하고 있는 경우에는 바로 컴파일에 성공하지 못할 수도 있다는 의미이다. 그 대신 \verb|mytblrpreamble.tex|이라는 이름의 파일을 활용하라고 하는 것이다.

예를 들면 다음과 같은 내용을 가진 \texttt{mytblrpreamble.tex} 파일을 작업 폴더에 만들어두고 표를 \texttt{article} 클래스 문서로 만들면 원하는 대로 컴파일이 가능해질 것이다. 

\begin{exampleonly}
\usepackage{margin=1in]{geometry}
\usepackage{kotex}
\input{svgnam.def}
\UseTblrLibrary{siunitx,booktabs}
%...
\end{exampleonly}

\noindent 이 예시에서 \verb|svgnam.def|에 대해 조금 말하자면 이것은 원래
\begin{verbatim}
\usepackage[svgnames]{xcolor}
\end{verbatim}
이렇게 하고 싶은 것인데 export되는 파일에 이미 \verb|\usepackage{xcolor}|가 들어 있어서 옵션 충돌이 일어난다. 이를 피하기 위하여 definition 파일만을 별도로 읽어들이게 한 것이다.

\href{https://nonametex.imweb.me/?bmode=view&idx=170507725&t=board}{noname의 ktug 임시 피난처} 게시판 글은 standalone tex export 기능과 \texttt{latexmk}의 \verb|-pvc| 옵션을 이용하여 마치 실시간 미리보기를 하듯이 작업할 수 있음을 보이고 있다.

\section{csv로 내보내기}

\tblrdoumi 로 작성하는 표 데이터를 Excel이나 스프레드시트로 보내는 방법으로 복사하여 붙여넣기가 가능함을 앞서 보였다.

그리고 \verb|csv| 포맷으로 내보낼 수 있다. 다만 괘선이나 셀 병합 정보는 유지되지 않는다. 데이터만 저장되는 것으로 생각하면 될 것이다. 앞서 작성했던 도표를 csv로 내보내기 하면 다음과 같은 결과를 얻게 된다. 

\begin{verbatim}
Value 1,Value 2,Value 3
$\alpha\beta$,,$gamma$
1234,,a
,,b
3,23.113231,c
,25.113231,d
\end{verbatim}

정말로 알맹이만 남기고 있음을 알 수 있다. 그래도 Excel이나 스프레드시트 프로그램으로 읽어들여서 편집할 수 있게 되므로 표의 내용을 전달하는 데는 문제가 없을 것이다.

\chapter{기타}

\section{\bs end 직전 명령}

기본 화면의 우하단 \guielem{LaTeX 코드 미리보기} 영역 아래쪽에 조그마하고 흐린 
체크박스가 하나 있고 ``\bs end 직전 명령''이라고 되어 있다. 이것은 무엇에 쓰는 것일까?

여기서 말하는 ``\bs end 직전''이란, 마지막 개행 명령(\verb|\\|)과 \verb|\end{tblr}|
사이를 말한다. 여기에 뭔가를 넣는데 왜 이런 것이 필요한지 생각해보자. 다음 \env{tabular} 코드를 보면,
\begin{exampleonly}
\begin{tabular}{lr}
\hline
A & B \\ \hline
1 & 2 \\ \hline
\end{tabular}
\end{exampleonly}
\verb|\hline| 명령을 맨 처음이나 \verb|\\| 직후에 이어서 쓰고 있다.

\tblrdoumi 는 `셀'을 파악할 때 \verb|&|와 \verb|\\|만을 기준으로 한다.
따라서 첫번째 \verb|\hline|은 위치상 (1,1) 셀에 들어 있는 것처럼 이해될 수밖에 없다.
이것은 뭔가 논리적이지 않기 때문에 괘선을 모두 ``설정부''에서 처리하는 것을 지향하는 것이다.
아무튼 그러면 마지막 \verb|\hline|은 어느 셀에 있는 것일까?

적어도 이 테이블에서 3행이라는 것은 존재하지 않는다. 따라서 (3,1)에 이 명령을 넣을 수는 없을 것이다. 
즉, \verb|&|와 \verb|\\|만으로 테이블을 분석할 때에 마지막 \verb|\hline|이 갈 곳이 없게 된다.

이 위치에 뭔가를 굳이 넣어야 한다면 ``\bs end 직전 명령''을 통해서 할 수 있다.

이것이 왜 필요했는가 하면 \pkg{tabularray}의 \pkg{booktabs} 라이브러리 때문이었다.
즉 마지막 \verb|\\| 이후에 \verb|\bottomrule|을 넣고 싶을 때를 위한 것.

\section{\LaTeX\ 코드 들여오기}

\hcrkey{F2} 단축키에 할당되어 있는 ``LaTeX 코드 불러오기''는 \env{tblr}, \env{longtblr}, \env{talltblr} 코드를 붙여넣기 하여 들여올 수 있다.

그리고 \pkg{tabularx} 패키지의 \env{tabularx} 환경도 (일부) 해석해서 가져온다.

\env{tabularx} 임포트 기능을 만들면서 가졌던 생각은 표의 알맹이를 가져오는 데 주력하고
장식적\.부수적인 정보는 웬만하면 없애버린다는 것이었다. 이와 관련하여 지적할 것이 있다.

\begin{itemize}\oblivoirlist
\item 괘선 정보는 모두 제거한다. 가로괘선, 부분괘선, 세로괘선 모두 없앤 상태로 들여온다. 괘선은 \tblrdoumi\ 내부에서 다시 그어야 한다.
\item 각 셀에 붙어 있는 (때로 엄청난) 부수적 매크로들은 정렬 관련, 폰트 관련 정보를 \cmdprint{\SetCell}로 변환해주려 하지만 복잡한 명령이라서 방법이 없으면 그대로 남겨둔다.
\item \env{tabularx}의 colspec은 \texttt{manual colspec}으로 보존한다.  \opt{manual} 라디오버튼을 선택하면 원래의 colspec을 편집하거나 적용할 수 있다.
\item 셀 병합 정보, 즉 \cmdprint{\multicolumn}과 \cmdprint{\multirow}는 가능한 한 해석하여 병합 정보로 활용할 수 있게 하려 했다. 특히 \cmdprint{\multicolumn} 때문에 어긋난 \verb|&| 개수를 채워넣는 데 주의를 기울였다.
\end{itemize}


\section{명령행 유틸리티}

\texttt{.tblr}이라는 파일 포맷을 정의함으로써 다른 프로그램과의 정보 교환이 수월해졌다.
다음에 소개하는 두 개의 명령행 유틸리티는 \tblrdoumi\ 기능의 일부를 독립 명령행 앱으로 만든
것이거나 일부 기능을 확장한 것이다.

\subsection{excel2tblr}

``Excel 불러오기'' 기능을 CLI 툴로 만든 것이다. 사용법은 다음과 같다.

\begin{center}
\includegraphics[width=\linewidth]{ffd0ccad6719e}
\end{center}

기본적으로 \verb|.tblr|로 저장한다. \verb|--tex| 옵션으로 \verb|.tex| 파일 저장이 가능하므로 어느 정도는 엑셀을 tex으로 변환하는 콘버터 역할을 할 수 있겠지만 결국 괘선이라든가 그런 세세한 설정을 재편집해야 할 필요가 있으므로 완벽한 \verb|xlsx2tex| 변환기라고는 할 수 없다.
이 역시 \tblrdoumi 에게 편집을 시작할 준비를 갖추게 하는 것이 목표이다. 특징이라면 Excel의 셀 병합 정보를 제법 잘 가져온다는 것과 숫자를 문자열 등과 다르게 처리하는 점이다.

\subsection{tabular2tblr}

\tblrdoumi 는 \env{tabularx}를 불러오는 기능을 가지고 있지만 현실적으로 \env{tabularx}보다는 \env{tabular}가 더 많을 것이다. 이 부분은 별도의 명령행 유틸리티로 제작하였다.

\env{tabular}는 40년 역사의 레거시이다. 수많은 변종과 변형이 있으므로 모든 \env{tabular}를 완벽하게 한방에 변환해주는 그런 것은 존재할 수 없다. 우리가 의도한 것은 단지 표의 내용을 잘 유지하고 \tblrdoumi 에서 추가적인 편집을 하기 쉽도록 바탕을 만드는 것이었다.

그냥 \verb|tabular2tblr|을 실행하면 ``붙여넣기 대기'' 상태가 된다. \env{tabular} 환경 부분을 복사한 다음 여기에 붙여넣고 \hcrkey{Enter}를 치고 \hcrkeys[+]{Ctrl,Z}를 누른다.
(Windows 터미널에서 ``여러 줄 붙여넣기 진짜 할 거냐?''고 물어오기도 하지만 그냥 진행하면 될 것이다. 붙여넣은 후에 \hcrkey{Enter}를 한번 눌러야 하는 것을 기억하자.)

별도로 저장한 파일이 있다면 \verb|-i| 옵션으로 파일 이름을 지정할 수 있는데 이 파일에는 \env{tabular} 이외의 내용이 들어 있으면 안된다.

출력 파일의 디폴트는 \verb|output.tblr|이지만 \verb|-o| 옵션으로 출력 파일 이름을 지정할 수 있다. 명령이 실행된 현재 폴더에 저장한다. 더블클릭하여 수정을 시작한다.

도움말을 보려면 \verb|-h| 옵션을 준다.

\section{권장 작업 순서}

이 앱으로 표를 그리려 할 때 작업 순서는 다음과 같다.
\begin{enumerate}[\quad 1.] \oblivoirlist
\item 표의 본체를 입력하거나 가져온다.
\item 행과 열, 셀의 속성을 설정한다.
\item 괘선을 그린다.
\item 셀 병합을 처리한다. 병합된 셀이 있으면 행\.열\.셀 속성이나 괘선을 그리기 위하여 셀을 선택할 때 불편할 수 있으므로 행\.열\.셀 속성과 괘선보다 나중에 처리하는 것이 좋다.
\item colspec을 정리하여 확정한다.
\end{enumerate}

이렇게 하여 완성된 표를 \hcrkey{F4}로 복사하거나 \hcrkeys[+]{Ctrl, T}로 내보내기 하여 문서에서 활용하는 것이다.





\end{document}
